% ReHwp-Bench-Paper-copy/kcc.tex의 공통 프리앰블에서 계승.
\usepackage{graphicx}
\usepackage{amsmath}
\usepackage{amssymb}
\usepackage{array}
\usepackage{booktabs}
\usepackage{tabularx}
\usepackage{float}
\usepackage{url}
\usepackage{caption}
\usepackage{tikz}
\usetikzlibrary{arrows.meta,positioning,patterns,decorations.pathreplacing}
\usepackage{pgfplots}
\pgfplotsset{compat=1.18}
\usepackage{balance}
% \raggedbottom에서도 각주를 단 하단에 배치한다.
\usepackage[bottom]{footmisc}

\captionsetup{font=small,skip=3pt,labelsep=colon}
\setlength{\abovecaptionskip}{3pt}
\setlength{\belowcaptionskip}{3pt}
\clubpenalty=10000
\widowpenalty=10000
% 짧은 2단 원고의 빈 공간을 절·문단 사이로 늘리지 않는다.
\raggedbottom
\renewcommand{\tablename}{표}
\renewcommand{\figurename}{그림}
\renewcommand{\refname}{참고문헌}
\newcommand{\ph}[1]{\ensuremath{\langle #1\rangle}}
\newcommand{\draftnote}[1]{\par\noindent\textbf{[작성 메모]} #1\par}
\newcolumntype{L}{>{\raggedright\arraybackslash}X}
\newcommand{\papertablestyle}{\small\renewcommand{\arraystretch}{1.15}}

\tikzset{
  paperbox/.style={draw,rounded corners=1.2pt,align=center,font=\footnotesize,
    inner xsep=1.6pt,inner ysep=1.5pt,minimum height=5mm,minimum width=26mm},
  paperresult/.style={paperbox,fill=black!7},
  paperarrow/.style={-{Latex[length=1.4mm,width=1mm]},semithick},
  paperlabel/.style={font=\scriptsize,fill=white,inner sep=0.7pt}
}
